\section{The exact remaining arithmetic comparison}\label{sec:remaining}
The fixed arithmetic presentation gives, with $a_0,a_1,a_2$ as in the polarized block,
\[
 C_{\rm pol}=\frac{\epsilon a_1a_2^3}{\omega},\qquad
 C_{\rm pol}U_{\rm alg}^3=R_+:=|\epsilon|^4|a_1|^6|a_2|^4>0.
\]
The modulus theorem gives $|C_{\rm pol}U_\Gamma^3|=R_+$. The exact statement still missing is
\begin{equation}\label{eq:remaining-reality}
 \Im(C_{\rm pol}U_\Gamma^3)=0.
\end{equation}
Reality together with a valid positive-real-part enclosure would give $C_{\rm pol}U_\Gamma^3=R_+$, hence $U_\Gamma^3=U_{\rm alg}^3$. A valid independent bound $|U_\Gamma/U_{\rm alg}-1|<1/2$ would select the root $U_\Gamma=U_{\rm alg}$. A numerical enclosure containing zero for the imaginary part proves no exact reality statement; equality remains unproved however small the displayed residual.

The new canonical product advances the scalar analytic comparison by supplying every divisor order, its normal convergence, a gamma-free anchor and the exact polynomial phase coefficients. Its branch-free cube \eqref{eq:cube} still contains an infinite transcendental product. Evaluating that regularized product arithmetically, or providing a phase-sensitive regulator/cocycle comparison, would be an additional theorem. No such theorem follows from the finite-theta obstruction, the modulus comparison, or the formal integer calculations.

The scope of the final results is therefore a fixed complex analytic reconstruction and a precise arithmetic boundary. There is no general Hilbert-12 theorem, no full gamma/class-field identification, and no priority assertion based on a literature search. The relative torus route of Section~\ref{sec:rigidity} is a sufficient conditional completion mechanism, with its missing hypotheses stated separately.
